\chapter{Mercadotecnia: la campaña ``El sur tiene espacio''}
\label{cap:s-campana}
\textit{Materia: Mercadotecnia digital. Pregunta: ¿a quién le habla la campaña, con qué palabras, por qué medio, cuándo y
cómo se sabe si funciona?}

\section{El método: segmentar, elegir y posicionar}
La campaña sigue los tres pasos clásicos de la mercadotecnia:
\begin{enumerate}
\item \textbf{Segmentar}: dividir a los posibles visitantes en grupos, con datos.
\item \textbf{Elegir}: a qué grupos hablarles.
\item \textbf{Posicionar}: qué lugar quiere ocupar la marca en la mente de esos grupos.
\end{enumerate}
La diferencia con una campaña común es que \textbf{cada paso sale de una cifra}, y cada rasgo dice si es un dato medido, un
cálculo o un supuesto.

\section{Lo que dicen los datos sobre quién viaja}
\begin{itemize}
\item \textbf{Quién visita el sur} (INAH 2025): en Oxtankah, el 94.8\,\% son mexicanos; en la Ruta, el 63.3\,\%.
\item \textbf{Quién llega al estado en avión}: 9,408,423 extranjeros a Cancún en 2025; 56.3\,\% de Estados Unidos y
  17.1\,\% de Canadá; 53.4\,\% mujeres.
\item \textbf{El dato que lo define todo}: al aeropuerto de Chetumal llegaron \textbf{250 extranjeros} en todo 2025. El
  extranjero no llega al sur en avión: entra por Cancún y baja por carretera, en Tren Maya, desde Belice o en crucero.
  Por eso el anuncio tiene que encontrarlo \emph{antes} de viajar o \emph{cuando ya está en Cancún}.
\item \textbf{Cómo se informan} (ENDUTIH 2025): el 92.1\,\% de la gente de Quintana Roo usa internet; de quienes tienen
  teléfono inteligente, 90.6\,\% usa mensajería y 80.4\,\% redes sociales.
\item \textbf{Qué los aleja y qué los atrae} (85,987 reseñas, capítulo~\ref{cap:s-patrones}): el ruido, la suciedad, el
  precio y las multitudes hunden una reseña; la calma y la cultura la protegen.
\end{itemize}

\section{Las dos viajeras ideales}
Una \emph{buyer persona} es el retrato de un cliente típico, construido con datos. Brandon eligió \textbf{dos}:

\begin{opciones}[La que vuelve al sur (nacional)]
\begin{itemize}
\item \textbf{Origen}: mexicana. Es el 94.8\,\% de quien visita Oxtankah \emph{(dato, INAH)}.
\item \textbf{Cuándo viaja}: de diciembre a abril; la campaña la invita en los meses con espacio \emph{(derivado de la forma
  del año)}.
\item \textbf{Qué valora}: calma y cultura \emph{(derivado de las reseñas)}.
\item \textbf{Qué la aleja}: ruido, suciedad, precio y multitudes \emph{(derivado)}.
\item \textbf{Cómo se informa}: mensajería y redes sociales \emph{(dato, ENDUTIH)}.
\item \textbf{De qué estado viene, su edad, su ingreso}: \emph{hueco}; no hay dato oficial por destino.
\end{itemize}
\end{opciones}

\begin{opciones}[La que baja del norte (Estados Unidos o Canadá)]
\begin{itemize}
\item \textbf{Origen}: de los 9.4 millones de extranjeros en Cancún, 56.3\,\% de Estados Unidos y 17.1\,\% de Canadá
  \emph{(dato)}.
\item \textbf{Sexo}: 53.4\,\% mujeres \emph{(dato)}.
\item \textbf{Por qué no llega sola}: solo 250 extranjeros aterrizaron en Chetumal \emph{(dato)}.
\item \textbf{Cuándo está}: más en marzo, menos en septiembre; Cancún en temporada alta de noviembre a abril \emph{(dato)}.
\item \textbf{Cómo se informa}: busca en Google y usa redes desde el destino \emph{(supuesto: no hay encuesta propia)}.
\end{itemize}
\end{opciones}

\begin{porque}
\textbf{Solo la nacional} dejaría fuera al extranjero que ya está a 334 km, en Cancún. \textbf{Solo la extranjera} ignoraría
que la Bahía es 95\,\% mexicana. Con dos, cada una recibe su mensaje y su canal.
\end{porque}

\section{La marca}
\begin{itemize}
\item \textbf{Nombre}: ``El sur tiene espacio'' (en inglés, ``The south has room''). Dice el dato (capacidad que sobra) y la
  promesa (sin multitudes) a la vez. Se descartaron ``Sur sin prisa'' (no dice dónde) y ``Del Caribe al sur'' (solo le
  habla a la extranjera).
\item \textbf{Lema}: ``Pirámides, bahía y sabor del sur, con espacio para disfrutarlos''.
\item \textbf{Personalidad}: Brandon pidió combinar los tres tonos que se le propusieron, ``porque es un orgullo ser
  mexicano'': cercana y tranquila, orgullosa de lo maya y lo mexicano, y con espíritu de aventura.
\item \textbf{Propuesta de valor}: el Caribe mexicano que todavía tiene espacio: pirámides con 15.5 veces menos visitantes
  que Tulum, una bahía tranquila y la comida del sur.
\item \textbf{Lo visual}: el sistema ``Sur mexicano'': rosa mexicano, amarillo cempasúchil, turquesa y añil en bloques, la
  greca maya y fotos reales comprobadas.
\end{itemize}

\section{Los anuncios: cada promesa con su dato}
Reglas: cada anuncio lleva el dato que lo respalda; \textbf{no se prometen precios ni tiempos de viaje} (no hay dato
oficial); las palabras salen de las reseñas de 5 estrellas; y el largo se revisa contra los límites de cada plataforma.

\begin{center}
\small
\begin{tabularx}{\linewidth}{llY}
\encabezado \h{Para} & \h{Dónde} & \h{Anuncio y su respaldo} \\
Nacional & Google & ``Pirámides mayas con espacio''. Tulum: 1,031,443 visitantes; la Ruta: 66,628. \\
\filagris Nacional & Facebook & ``La bahía que te esperaba''. Oxtankah: 11,017 visitantes en 2025, 94.8\,\% mexicanos. \\
Nacional & WhatsApp & ``¿Y si este puente nos vamos al sur?''. 90.6\,\% usa mensajería; no cuesta. \\
\filagris Extranjera & Google (inglés) & ``Mayan pyramids, with room''. El mismo 15.5; Cancún en temporada alta. \\
Extranjera & Facebook (inglés) & ``Your Caribbean trip, one stop south''. Solo se muestra cuando la Torre marca el norte
lleno. \\
\bottomrule
\end{tabularx}
\normalsize
\end{center}

\section{Los medios y el calendario}
\begin{itemize}
\item \textbf{Facebook e Instagram, 70\,\%}: trae 6.61 visitantes por cada \$1,000 y el 80.4\,\% usa redes. Con tope para no
  depender de una sola plataforma.
\item \textbf{Google, 30\,\%}: llega a quien ya está buscando viajar.
\item \textbf{WhatsApp, sin costo}: el botón de compartir viaja entre familia y amigos.
\item \textbf{La página}: el destino de todos los anuncios, con el planeador por lugar y mes en 9 idiomas (el décimo, el maya yucateco, espera la revisión
  de una persona hablante antes de publicarse).
\end{itemize}
\textbf{Calendario}: los 9 meses del plan (octubre de 2026 a junio de 2027), con el dinero de la Fase 6. ``La que baja del
norte'' se activa solo en los meses en que Cancún está en temporada alta y el sur tiene espacio, y además la Torre la
enciende por semana cuando el norte se llena.

\section{Cómo se sabe si funciona: los indicadores}
\begin{center}
\small
\begin{tabularx}{\linewidth}{lYY}
\encabezado \h{Indicador} & \h{Cómo se calcula} & \h{Meta} \\
Alcance & Personas que vieron un anuncio & Se fija con el primer mes real \\
\filagris Interacción & Clics entre veces que se vio & Google 8.73\,\% o más; Facebook 2.76\,\% o más \\
Conversión & Viajes planeados entre clics & 5.75\,\% o más \\
\filagris Costo por visitante & Pesos gastados entre visitantes & \$196 o menos \\
Afluencia & Visitantes reales menos el escenario probable & 957 más en 9 meses \\
\filagris Presupuesto & Pesos gastados entre pesos planeados & 100\,\% (lo pausado se gasta después) \\
Económico & Ocupación hotelera de Chetumal & \textbf{Hueco}: no se publica desde 2025 \\
\filagris Social & Visitantes por cada mil habitantes (Radar) & Nunca ``saturado'' \\
Ambiental & Semanas con anuncio en temporada alta o con mal clima & 0 \\
\filagris Capacidad & Esperados más campaña entre capacidad probada & 40\,\% o menos \\
\bottomrule
\end{tabularx}
\normalsize
\end{center}

\section{La regla anti-colapso}
El incidente de mercadotecnia pregunta qué hacer si llega más gente de la que un lugar soporta. La campaña tiene la
respuesta \textbf{programada}, no solo escrita:
\begin{enumerate}
\item Nunca anuncia un lugar en su temporada alta (regla del presupuesto).
\item No permite que los visitantes esperados más los de la campaña pasen de la capacidad probada.
\item La Torre pausa el anuncio si hay tormenta, clima raro o llegó más gente de la esperada.
\item A quien iba al norte lleno le ofrece el lugar del sur con más espacio, nunca otro lugar lleno.
\end{enumerate}
